\section{Introduction}

Can a model's weight changes tell us how its measured capabilities will change, before we run the edited model? This question matters when evaluating a large collection of fine-tunes or interventions is expensive. Weight-space methods offer a possible alternative: instead of repeatedly executing every candidate, a predictor could use the parameters themselves to estimate the direction and approximate size of performance changes. Such estimates would not replace evaluation. They could help identify which edits merit a full benchmark run and expose capability costs that are easy to miss when only the intended behavior is measured.

The immediate motivation is the report by Takehi et al. that editing 0.016\% of Qwen3.5-0.8B's weights raised its MATH-500 score from 7.9\% to 48.1\% without gradients \citep{takehi2026wsp}. This is a striking instance of a weight edit with a measured behavioral effect, not evidence that the effect can be read from a delta in advance or that the result generalizes to other edit families. We ask a narrower, testable question: given a base model, a sequence of weight edits, and no new forward passes, can a predictor estimate the edited model's change on a fixed battery of benchmarks? We treat each transition as a delta $\Delta\theta_t=\theta_t-\theta_{t-1}$ and predict a vector of benchmark changes, rather than a single score or a natural-language account of the edit.

Our study is organized around four hypotheses. H1 asks whether weight-only representations predict the sign and rough magnitude of per-benchmark changes on held-out edits to the same base model. The planned criterion is a per-benchmark Spearman correlation of at least $0.6$, with the predictor compared against a statistics-based model and a delta-norm baseline. H2 asks whether modeling edit order improves prediction over treating each edit independently. We will compare sequence predictors with a stateless predictor on trajectories of at least four edits; the target is at least a 15\% relative reduction in mean absolute error. H3 tests transfer across edit families: a JEPA-style encoder trained on abliteration and random-direction edits will be evaluated on LoRA fine-tunes it has not seen, with a target rank correlation of at least $0.5$. These thresholds are planned decision criteria, not results. H4 compares an LLM asked to read the same weight-delta summaries with a trained predictor. We expect the trained predictor to be more accurate on familiar edit families and ask whether the LLM reader is more robust to novel families; both directions will be reported.

An alignment-edit sub-study will make the prediction problem concrete across model scales. For Qwen3.5-0.8B, 4B, and 9B and Qwen3.8-27B, we plan to compare base-instruct checkpoints, classic directional ablations, and Heretic edits. Evaluation will include capability benchmarks, refusal and over-refusal rates, and KL divergence from the base. We will then test whether the predictor can identify which edit family produced a delta and estimate associated capability changes. The measures are complementary: a refusal marker is not a complete measure of harmfulness, and a small KL divergence does not establish that task accuracy is preserved. The comparison is an audit of existing edits, not a proposal for a new abliteration method.

The planned work has three contributions. First, we will construct a controlled set of single edits and multi-edit trajectories with paired benchmark outcomes, so that prediction can be evaluated under splits by edit family and seed. Second, we will compare inexpensive weight statistics and canonicalized low-rank features with stateless, sequential, and JEPA-style predictors, alongside an LLM reader. This directly tests whether sequence context or a learned representation adds value beyond per-edit summaries. Third, we will apply the same prediction framing to alignment edits and report capability, refusal, and KL measures together. The central limitation is that this is a study of a planned edit zoo, initially centered on Qwen models; conclusions will depend on its edit-family coverage and evaluation battery. We do not yet report experimental outcomes, and we do not claim that weight-only prediction can replace benchmark evaluation.
